% GENERATED FILE. Edit canonical lessons, then run npm run generate:books.
% canonical-source-hash: fnv1a64:ab65f203d9845426
% canonical-lessons: TA-C404-vakanam, TA-C404-lari, TA-C404-skuttar, TA-C404-muntu, TA-C404-motu

\chapter{Lorries and Scooters}
\label{ch:ta-lorries-and-scooters}
\hlchaptermodality{404}

\begin{chapteropening}
\textbf{By the end of this chapter:} \emph{I can talk about vehicles, lorries and scooters, and overtaking and crashing.}

Complete the last lesson of chapter 404: I can talk about vehicles, lorries and scooters, and overtaking and crashing.
\end{chapteropening}

\section[\texorpdfstring{\ta{வாகனம்} (vākaṉam)}{வாகனம் (vākaṉam)}]{\ta{வாகனம்} (vākaṉam) — a motor vehicle; a vehicle}

\label{lesson:TA-C404-vakanam}



\begin{quote}
Before the new one: say the Tamil for a junction, then the Tamil for to speed.
\end{quote}

\subsection*{You'll want to know: \ta{வாகனம்}}
\textbf{\ta{வாகனம்}} — \emph{vākaṉam} — \textquotedblleft{}a motor vehicle; a vehicle\textquotedblright{}.

\begin{quote}
\ta{இந்தச்} \ta{சாலையில்} \ta{பெரிய} \ta{வாகனங்கள்} \ta{போகக்} \ta{கூடாது}.
\end{quote}

\begin{quote}
\emph{intac cālaiyil periya vākaṉaṅkaḷ pōkak kūṭātu.}
\end{quote}

\begin{quote}
Large vehicles must not go on this road.
\end{quote}

Plural: \ta{வாகனங்கள்}, \emph{vākaṉaṅkaḷ}.

A -\ta{ம்} noun makes its plural in -\ta{ங்கள்}: \ta{வாகனங்கள்}.

\begin{cousinweb}
\textbf{\ta{வாகனம்}} (\emph{vākaṉam}) is borrowed from Sanskrit \emph{vāhana}, \textquotedblleft{}a vehicle, a riding or draught animal\textquotedblright{} (CDIAL 11610); Tamil adds \emph{-m} to the final \emph{-a}.
\end{cousinweb}

\begin{grammarlens}[title={What you've built}]
One more word for getting around.
\end{grammarlens}

\subsection*{Your turn}
\begin{itemize}
\raggedright
  \item \emph{Say it:} \emph{vākaṉam}
  \item \emph{Say it:} \emph{vākaṉam}, once more
  \item \emph{Say it:} \emph{intac cālaiyil periya vākaṉaṅkaḷ pōkak kūṭātu.} — Large vehicles must not go on this road.
  \item \emph{Recall:} say the Tamil for a junction, then the Tamil for to speed, then say \emph{vākaṉam} again
\end{itemize}

\begin{tcolorbox}[breakable,colback=teal!4,colframe=teal!35!black,title={Before you move on}]
What does \ta{வாகனம்} mean? (\textquotedblleft{}A motor vehicle; a vehicle\textquotedblright{}.) Say it once more.
\end{tcolorbox}
\section[\texorpdfstring{\ta{லாரி} (lāri)}{லாரி (lāri)}]{\ta{லாரி} (lāri) — a lorry; a truck}

\label{lesson:TA-C404-lari}



\begin{quote}
Before the new one: say the Tamil for to speed, then the Tamil for a motor vehicle.
\end{quote}

\subsection*{You'll want to know: \ta{லாரி}}
\textbf{\ta{லாரி}} — \emph{lāri} — \textquotedblleft{}a lorry; a truck\textquotedblright{}.

\begin{quote}
\ta{ஒரு} \ta{லாரி} \ta{மணல்} \ta{கொண்டு} \ta{வந்தது}.
\end{quote}

\begin{quote}
\emph{oru lāri maṇal koṇṭu vantatu.}
\end{quote}

\begin{quote}
A lorry brought sand.
\end{quote}

The plural is \ta{லாரிகள்}.

\begin{cousinweb}
\textbf{\ta{லாரி}} (\emph{lāri}) is a loan from English \emph{lorry}.
\end{cousinweb}

\begin{grammarlens}[title={What you've built}]
One more word for getting around.
\end{grammarlens}

\subsection*{Your turn}
\begin{itemize}
\raggedright
  \item \emph{Say it:} \emph{lāri}
  \item \emph{Say it:} \emph{lāri}, once more
  \item \emph{Say it:} \emph{oru lāri maṇal koṇṭu vantatu.} — A lorry brought sand.
  \item \emph{Recall:} say the Tamil for to speed, then the Tamil for a motor vehicle, then say \emph{lāri} again
\end{itemize}

\begin{tcolorbox}[breakable,colback=teal!4,colframe=teal!35!black,title={Before you move on}]
What does \ta{லாரி} mean? (\textquotedblleft{}A lorry; a truck\textquotedblright{}.) Say it once more.
\end{tcolorbox}
\section[\texorpdfstring{\ta{ஸ்கூட்டர்} (skūṭṭar)}{ஸ்கூட்டர் (skūṭṭar)}]{\ta{ஸ்கூட்டர்} (skūṭṭar) — a scooter}

\label{lesson:TA-C404-skuttar}



\begin{quote}
Before the new one: say the Tamil for a motor vehicle, then the Tamil for a lorry.
\end{quote}

\subsection*{You'll want to know: \ta{ஸ்கூட்டர்}}
\textbf{\ta{ஸ்கூட்டர்}} — \emph{skūṭṭar} — \textquotedblleft{}a scooter\textquotedblright{}.

\begin{quote}
\ta{அக்கா} \ta{ஸ்கூட்டரில்} \ta{அலுவலகத்துக்குப்} \ta{போகிறாள்}.
\end{quote}

\begin{quote}
\emph{akkā skūṭṭaril aluvalakattukkup pōkiṟāḷ.}
\end{quote}

\begin{quote}
My elder sister goes to the office on a scooter.
\end{quote}

\ta{ஸ்} is the grantha s, kept in this loan.

\begin{cousinweb}
\textbf{\ta{ஸ்கூட்டர்}} (\emph{skūṭṭar}) is a loan from English \emph{scooter}.
\end{cousinweb}

\begin{grammarlens}[title={What you've built}]
One more word for getting around.
\end{grammarlens}

\subsection*{Your turn}
\begin{itemize}
\raggedright
  \item \emph{Say it:} \emph{skūṭṭar}
  \item \emph{Say it:} \emph{skūṭṭar}, once more
  \item \emph{Say it:} \emph{akkā skūṭṭaril aluvalakattukkup pōkiṟāḷ.} — My elder sister goes to the office on a scooter.
  \item \emph{Recall:} say the Tamil for a motor vehicle, then the Tamil for a lorry, then say \emph{skūṭṭar} again
\end{itemize}

\begin{tcolorbox}[breakable,colback=teal!4,colframe=teal!35!black,title={Before you move on}]
What does \ta{ஸ்கூட்டர்} mean? (\textquotedblleft{}A scooter\textquotedblright{}.) Say it once more.
\end{tcolorbox}
\section[\texorpdfstring{\ta{முந்து} (muntu)}{முந்து (muntu)}]{\ta{முந்து} (muntu) — to overtake; to get ahead}

\label{lesson:TA-C404-muntu}



\begin{quote}
Before the new one: say the Tamil for a lorry, then the Tamil for a scooter.
\end{quote}

\subsection*{You'll want to know: \ta{முந்து}}
\textbf{\ta{முந்து}} — \emph{muntu} — \textquotedblleft{}to overtake; to get ahead\textquotedblright{}.

\begin{quote}
\ta{அந்த} \ta{லாரி} \ta{எங்கள்} \ta{காரை} \ta{முந்தியது}.
\end{quote}

\begin{quote}
\emph{anta lāri eṅkaḷ kārai muntiyatu.}
\end{quote}

\begin{quote}
That lorry overtook our car.
\end{quote}

Past (I …): \ta{முந்தினேன்}, \emph{muntiṉēṉ}.

What is overtaken takes -\ta{ஐ}: \ta{காரை} \ta{முந்து}, to overtake a car.

\begin{cousinweb}
\textbf{\ta{முந்து}} (\emph{muntu}) is native Dravidian (DEDR 5020), the family of \ta{முன்}, \textquotedblleft{}front\textquotedblright{}.
\end{cousinweb}

\begin{grammarlens}[title={What you've built}]
One more word for getting around.
\end{grammarlens}

\subsection*{Your turn}
\begin{itemize}
\raggedright
  \item \emph{Say it:} \emph{muntu}
  \item \emph{Say it:} \emph{muntu}, once more
  \item \emph{Say it:} \emph{anta lāri eṅkaḷ kārai muntiyatu.} — That lorry overtook our car.
  \item \emph{Recall:} say the Tamil for a lorry, then the Tamil for a scooter, then say \emph{muntu} again
\end{itemize}

\begin{tcolorbox}[breakable,colback=teal!4,colframe=teal!35!black,title={Before you move on}]
What does \ta{முந்து} mean? (\textquotedblleft{}To overtake; to get ahead\textquotedblright{}.) Say it once more.
\end{tcolorbox}
\section[\texorpdfstring{\ta{மோது} (mōtu)}{மோது (mōtu)}]{\ta{மோது} (mōtu) — to collide; to crash into}

\label{lesson:TA-C404-motu}



\begin{quote}
Before the new one: say the Tamil for a scooter, then the Tamil for to overtake.
\end{quote}

\subsection*{You'll want to know: \ta{மோது}}
\textbf{\ta{மோது}} — \emph{mōtu} — \textquotedblleft{}to collide; to crash into\textquotedblright{}.

\begin{quote}
\ta{ஒரு} \ta{ஸ்கூட்டர்} \ta{மரத்தில்} \ta{மோதியது}.
\end{quote}

\begin{quote}
\emph{oru skūṭṭar marattil mōtiyatu.}
\end{quote}

\begin{quote}
A scooter crashed into a tree.
\end{quote}

Past (I …): \ta{மோதினேன்}, \emph{mōtiṉēṉ}.

What is hit takes -\ta{இல்} or -\ta{உடன்}: \ta{மரத்தில்} \ta{மோது}, to crash into a tree.

\begin{cousinweb}
\textbf{\ta{மோது}} (\emph{mōtu}) is native Dravidian (DEDR 5117), \textquotedblleft{}to strike, to dash against\textquotedblright{}.
\end{cousinweb}

\begin{grammarlens}[title={What you've built}]
One more word for getting around.
\end{grammarlens}

\subsection*{Your turn}
\begin{itemize}
\raggedright
  \item \emph{Say it:} \emph{mōtu}
  \item \emph{Say it:} \emph{mōtu}, once more
  \item \emph{Say it:} \emph{oru skūṭṭar marattil mōtiyatu.} — A scooter crashed into a tree.
  \item \emph{Recall:} say the Tamil for a scooter, then the Tamil for to overtake, then say \emph{mōtu} again
\end{itemize}

\begin{tcolorbox}[breakable,colback=teal!4,colframe=teal!35!black,title={Before you move on}]
What does \ta{மோது} mean? (\textquotedblleft{}To collide; to crash into\textquotedblright{}.) Say it once more.
\end{tcolorbox}
